\bigskip
\begin{center}\bfseries
II. Traditionen über die erste Jugendzeit bis zum Abgang auf die Universität\par
Oktober 1795.
\end{center}
\indent \textsc{Gauss}’ Mutter, die als eine feinsinnige Frau von festem und heiterem
Charakter geschildert wird, ist 1817 als 74jährige Witwe zu ihrem Sohne
nach Göttingen gezogen, wo sie noch 22 Jahre bei ihm auf der Sternwarte\par
\HistoricalFoot{X\textsubscript{2} Abh. 2.}{2}
\clearpage\SourceLeaf{95}\ArithmeticHead{10}{L}{L. SCHLESINGER, ÜBER GAUSS’ FUNKTIONENTHEORETISCHE ARBEITEN.}
\noindent gewohnt hat. Ihren Erzählungen verdankt man Nachrichten über die geistige
Entwicklung von \textsc{Gauss} in den ersten Jahren seiner Kindheit\footnote{Vergl. für die im folgenden wiedergegebenen Einzelheiten: \textsc{Sartorius v. Waltershausen}, \textit{Gauss zum Gedächtnis}, Leipzig 1856, S. 12ff.; L. \textsc{Hänselmann}, \textit{Carl Friedrich Gauss}, Leipzig 1878, S. 16ff.}.

Schon kaum dreijährig habe er den Vater bei der Lohnauszahlung am
Samstag Abend auf einen Rechenfehler aufmerksam gemacht, das Lesen habe
er sich selbst ohne Unterricht angeeignet. Aus seinem siebenten Lebensjahr
erzählt \textsc{Gauss} selbst, er habe in der Schule eine langwierige Rechenaufgabe
in wenigen Augenblicken gelöst, indem er auf das Gesetz der arithmetischen
Reihe aufmerksam wurde, und während seine Kameraden sich eben anschickten,
durch mühsames Addieren die Lösung zu finden, habe er seine nur eine einzige
Zahl enthaltende Tafel mit dem übermütigen Ausruf im Braunschweiger
Dialekt: »Dar licht se« auf das Katheder geworfen. Der Eindruck, den diese
Leistung auf den Schulmeister \textsc{Büttner} hervorgerufen hat, bewirkte, dass
dieser für den ungewöhnlichen Schüler aus Hamburg ein Rechenbuch kommen
liess, \textsc{Remers} \textit{Arithmetica}, das sich mit der Eintragung: »\textsc{Johann Friedrich
Carl Gauss}, Braunschweig, 16. December Anno 1785« noch in der Gaussbibliothek
befindet und ebenso wie das Exemplar von \textsc{Hemelings} \textit{Arithmetischem
kleinen Rechenbuch} Spuren starker Benutzung und zwischen dem Text einige
von \textsc{Gauss}’ kindlicher Hand ausgeführte elementare Rechnungen zeigt. — Von
grosser Bedeutung war für \textsc{Gauss} der Einfluss, den ein nur wenige Jahre
älterer Hilfslehrer an der \textsc{Büttner}schen Schule \textsc{Johann Martin Christian
Bartels} (geb. 1769) auf ihn ausübte. — \textsc{Bartels}, der später Professor der
Mathematik an den Universitäten Kasan und Dorpat war, hat \textsc{Gauss} nicht
nur in die alten Sprachen eingeführt, sondern ihn auch mit dem binomischen
Lehrsatz und der Lehre von den unendlichen Reihen bekannt gemacht, und
seinem Einfluss ist es wohl zu danken, dass \textsc{Gauss} 1788 in das Gymnasium
Katharineum eintreten durfte. \textsc{Bartels} führte, nachdem er selbst das Collegium
Carolinum (aus dem die jetzige Technische Hochschule hervorgegangen
ist) bezogen hatte, \textsc{Gauss} zu dem Professor der Mathematik an dieser Anstalt,
\textsc{August Wilhelm Zimmermann}, der eine warme Zuneigung zu dem begabten
Knaben fasste. Als 14jähriger Primaner wurde \textsc{Gauss} 1791 zum ersten Male
beim herzoglichen Hofe vorgestellt, wo er seine Rechenkünste zeigen durfte;
von da ab hat er unter dem besonderen Schutze des Herzogs \textsc{Carl Wilhelm}\par
\clearpage\SourceLeaf{96}\ArithmeticHead{11}{R}{II. DIE ERSTE JUGENDZEIT BIS OKTOBER 1795.}
\noindent \textsc{Ferdinand} gestanden, der sein lebelang in grosszügiger Weise für \textsc{Gauss} gesorgt
hat. Bei Gelegenheit der ersten Vorstellung bei Hofe erhielt \textsc{Gauss}
von dem Geheimrat \textsc{Feronce von Rothenkreuz} einige mathematische Lehrbücher
und Tafelwerke zum Geschenk, darunter die bekannten Tafeln von
\textsc{Schulze}\footnote{\textsc{Johann Carl Schulze}, \textit{Neue und erweiterte Sammlung logarithmischer .\ .\ . Tafeln} I, II, Berlin 1778.}, die \textsc{Gauss} auch später noch viel benutzt hat und in die er seinen
Namen mit der Jahreszahl 1791 eingetragen hat. So ist es nicht verwunderlich,
dass ihm das Jahr 1791 in dauernder Erinnerung geblieben ist.

In einem April 1816 datierten Briefe an \textsc{Schumacher}\footnote{Werke X, 1, S. 247; wir kommen auf diesen Brief noch zurück.} schreibt \textsc{Gauss}:
»Haben Sie denn wirklich vergessen, dass das arithmetisch-geometrische Mittel,
mit welchem Hr. \textsc{Degen} sich beschäftigt, \textls[120]{ganz dasselbe} ist, womit ich mich
seit 1791 beschäftigt habe, und jetzt einen ziemlichen Quartband darüber
schreiben könnte?« Hiernach hat also \textsc{Gauss} im Jahre 1791 begonnen, sich
mit dem agM. zu beschäftigen\footnote{Die Abkürzung agM. für arithmetisch-geometrisches Mittel wird auch im folgenden stets angewandt.}.

Fast in allen Fällen, wo \textsc{Gauss} im reiferen Alter in Briefen oder sonst
wo eine weit zurückliegende Zeitangabe für ein auf sein Leben oder seine
wissenschaftliche Entwicklung bezügliches Ereignis macht, lässt sich nachweisen,
dass er sich dabei nicht auf sein, übrigens hervorragend treues Gedächtnis
verlassen hat, sondern in irgend einer Aufzeichnung einen untrüglichen Anhaltspunkt
besass. Man wird also annehmen können, dass der zu \textsc{Schumacher}
getanen Äusserung auch etwas dokumentarisches zugrunde lag, und es liegt
nahe hier an den \textit{Schulze} zu denken, zumal \textsc{Gauss} auf unbedruckten Blättern
dieses Buches eine Reihe von Ergebnissen eingetragen hat, die sich auf die
Bestimmung mittlerer Werte der Zahlentheorie beziehen (vergl. Werke X, 1,
S. 11) und, wie sich auch später noch einmal zeigen wird, diesen an sich so
verschiedenartigen Mittelbildungen, wenigstens in der frühen Zeit, ein gemeinsames
Interesse als Ausgangspunkt gedient hat. — \textsc{Gauss} dürfte in jener
Jünglingszeit, als seine Vorliebe für mathematische Gegenstände, Zahlenverhältnisse
und ihre Anwendungen anfing festere Formen anzunehmen, überhaupt
auf die Bedeutung der Mittelbildung aufmerksam geworden sein. — Er mag
nun etwa für zwei Zahlen, z. B. zwei Näherungswerte einer zwischen ihnen
liegenden Grösse, zuerst das arithmetische Mittel gebildet, dann mit Hilfe\par
\HistoricalFoot{}{2\textsuperscript{*}}
\clearpage\SourceLeaf{97}\ArithmeticHead{12}{L}{L. SCHLESINGER, ÜBER GAUSS’ FUNKTIONENTHEORETISCHE ARBEITEN.}
\noindent des \textit{Schulze}, seiner ersten Logarithmentafel, das arithmetische Mittel der Logarithmen
dieser Zahlen ausgerechnet und den zugehörigen Numerus, also das
geometrische Mittel, aufgesucht haben. Er bemerkt dann, dass die Differenz
der beiden Mittel kleiner ist, als die der ursprünglichen Zahlen und wiederholt
darum das Verfahren. — Vielleicht hat er sich auch schon eine Vorstellung
von dem Grenzwert zu bilden gewusst, dem die Folgen dieser Mittel
sich sehr rasch annähern. — Dass \textsc{Gauss} in jenen Jahren die Mittelbildung
im Zusammenhang mit Aufgaben der Fehlerausgleichung zum Gegenstande
seiner Überlegungen gemacht hat, geht schon daraus hervor, dass er nicht
lange nachher zu dem Prinzip der Methode der kleinsten Quadrate vorgedrungen
sein muss, dessen er sich nach seinen Angaben seit 1794\footnote{So in Briefen an \textsc{Olbers} vom 30. Juni 1806 und an \textsc{Schumacher} vom 3. Dezember 1831, Werke VIII, S. 138, 139, 141. Diesen aus weit auseinanderliegenden Zeiten stammenden übereinstimmenden Zeugnissen gegenüber muss die Angabe in der \textit{Theoria Motus} (1809, Werke VII, S. 253), die das Jahr 1795 als dasjenige nennt, von dem ab \textsc{Gauss} sich jenes Prinzips bedient haben will, als irrtümlich angesehen werden; vergl. auch Werke X, 1, S. 373, 380 und XII, S. 64ff.} bedient hat.

Die eben entwickelten Vermutungen über den Weg, der \textsc{Gauss} zum agM.
geführt haben mag, könnten überflüssig erscheinen, weil ja \textsc{Lagrange} schon
1784—85\footnote{Siehe \textsc{Lagrange}, \textit{Sur une nouvelle méthode de calcul intégral pour les différentielles affectées d’un radical sous lequel la variable ne passe pas le quatrième degré}, Mémoires de Turin 2, 1784—85, S. 237, Oeuvres II, 1868, S. 251, siehe insbesondere S. 264.} den Algorithmus des agM. aufgestellt hat. — \textsc{Lagrange} geht dort
von der \textsc{Landen}schen Transformation
\[
y'=\frac{y\sqrt{(1\pm p^2y^2)(1\pm q^2y^2)}}{1\pm q^2y^2}
\]
aus, durch die sich
\[
\frac{dy}{\sqrt{(1\pm p^2y^2)(1\pm q^2y^2)}}\quad\text{in}\quad
\frac{dy'}{\sqrt{(1\pm p'^2y'^2)(1\pm q'^2y'^2)}}
\]
verwandelt, wo
\[
p'=p+\sqrt{p^2-q^2},\qquad q'=p-\sqrt{p^2-q^2}
\]
ist; für
\[
p+q=m,\quad p-q=n,\quad p'+q'=m',\quad p'-q'=n'
\]
hat man also
\[
m'=m+n,\qquad n'=2\sqrt{mn}.
\]

\clearpage\SourceLeaf{98}\ArithmeticHead{13}{R}{II. DIE ERSTE JUGENDZEIT BIS OKTOBER 1795.}
\noindent Er bemerkt dann (a. a. O. S. 271), dass bei wiederholter Anwendung dieser
Transformation, die sich ergebenden Folgen \(p,p',p'',\ldots\) und \(q,q',q'',\ldots\) von
einander weg divergieren (sont divergentes l’une par rapport à l’autre), sodass
man also, indem man sie nach rückwärts hin fortsetzt, zwei Folgen \(\ldots\,{}''p,{}'p,p\)
und \(\ldots\,{}''q,{}'q,q\) erhält, die gegeneinander konvergieren. Es ist dann
\[
p={}'p+\sqrt{{}'p^2-{}'q^2},\quad q={}'p-\sqrt{{}'p^2-{}'q^2}\quad\text{usw.},
\]
also
\[
{}'p=\frac{p+q}{2},\quad {}'q=\sqrt{pq};\quad
{}''p=\frac{{}'p+{}'q}{2},\quad {}''q=\sqrt{{}'p{}'q};\quad\text{usw.},
\]
sodass die korrespondierenden Glieder stets das arithmetische bezw. geometrische
Mittel der vorhergehenden sind. Wenn nun \(p>q\), so hat man offenbar
\({}'p<p,\ {}'q>q,\ {}'q<{}'p\), weil
\[
{}'p-{}'q=\frac{p+q}{2}-\sqrt{pq}=\tfrac12(\sqrt p-\sqrt q)^2;
\]
die Reihe \(p,{}'p,{}''p,\ldots\) ist also eine abnehmende, die Reihe \(q,{}'q,{}''q,\ldots\) dagegen
eine zunehmende, und der Unterschied der korrespondierenden Glieder
nimmt ins Unendliche ab. Setzt man dann wieder
\[
y=\frac{{}'y\sqrt{(1\pm{}'p^2{}'y^2)(1\pm{}'q^2{}'y^2)}}{1\pm{}'q^2{}'y^2}\quad\text{usw.}
\]
so ist
\[
\frac{dy}{\sqrt{(1\pm p^2y^2)(1\pm q^2y^2)}}=
\frac{d{}'y}{\sqrt{(1\pm{}'p^2{}'y^2)(1\pm{}'q^2{}'y^2)}}=\text{usw.},
\]
Dieser Algorithmus wird zur näherungsweisen Berechnung des Integrals
\[
\int\frac{M\,dy}{\sqrt{(1\pm p^2y^2)(1\pm q^2y^2)}}\mathpunct{\vcenter{\hbox{,}}}
\]
wo \(M\) eine beliebige rationale Funktion von \(y^2\) bedeutet, und (S. 283) insbesondere
auf das bei der Rektifikation der Ellipse und Hyperbel auftretende
Integral
\[
\int\frac{(1-e^2x^2)\,dx}{\sqrt{(1-x^2)(1-e^2x^2)}}
\]
angewandt. Man nennt den gemeinsamen Grenzwert, dem die Folgen \(p,{}'p,
{}''p,\ldots\) und \(q,{}'q,{}''q,\ldots\) zustreben, nach \textsc{Gauss}, das agM. zwischen \(p\) und \(q\)
und bezeichnet es mit \(M(p,q)\).

Wenn es schon an sich sehr unwahrscheinlich ist dass der 14jährige
Gymnasiast \textsc{Lagrange} gelesen haben sollte, so besitzen wir noch eine aus-\par
\clearpage\SourceLeaf{99}\ArithmeticHead{14}{L}{L. SCHLESINGER, ÜBER GAUSS’ FUNKTIONENTHEORETISCHE ARBEITEN.}
\noindent drückliche Erklärung von \textsc{Gauss}, der in der Anzeige der Abhandlung \textit{Determinatio attractionis etc.}
(1818, Werke III, S. 360) sagt, dass er die auf das
agM. bezüglichen Resultate, die er in jener Abhandlung zum ersten Male
von sich aus veröffentlicht, »so wie er sie schon vor vielen Jahren unabhängig
von ähnlichen Untersuchungen \textsc{Lagranges} und \textsc{Legendres} gefunden
hat, in ihrer ursprünglichen Form darstellen zu müssen geglaubt hat, .\ .\ . teils
weil jene Form ihm wesentliche Vorzüge zu haben schien, teils weil sie gerade
so den Anfang einer viel ausgedehnteren Theorie ausmachen, wo seine
Arbeit eine ganz verschiedene Richtung von der der genannten Geometer genommen
hat«. Freilich konnte der wissbegierige Knabe bei seiner mehr
spielerischen Beschäftigung mit dem agM. nicht ahnen, dass er sich mit diesem
Problem eine schier unerschöpfliche Quelle reichster und tiefster Erkenntnisse
erbohrt hatte, aus der noch der reife Mann bis ins hohe Alter immer wieder
schöpfen konnte, aber es scheint doch, als ob ihn ein unbewusstes Gefühl
für die eigenartige Bedeutung des agM. veranlasst hätte, den Gegenstand nicht
wieder aus dem Auge zu lassen. Allerdings besitzen wir hierfür aus den
nächsten Jahren keine Zeugnisse, wir wissen nur, dass \textsc{Gauss}, als er 1792
vom Gymnasium abgegangen war, einer bestehenden Verordnung gemäss, ehe
er die Universität beziehen durfte, erst noch den dreijährigen Kurs des Collegium
Carolinum absolvieren musste und dass er sich am 18. Februar 1792
in die Matrikel dieses Instituts eingeschrieben hat. Wie der bereits erwähnte
Professor \textsc{Zimmermann} im Jahre 1796 bezeugt\footnote{Siehe Werke X, 1, S. 3.}, hat \textsc{Gauss} sich auf dem Carolinum
»mit ebenso glücklichem Erfolge der Philosophie und der klassischen
Literatur als der höheren Mathematik gewidmet.« Erst für das Jahr 1794
berichtet \textsc{Schering} (Werke III, S. 493) »nach Mitteilungen über eine mündliche
Äusserung von \textsc{Gauss}«, dass dieser 1794 »die Beziehungen zwischen dem
agM. und den Potenzreihen, in denen die Exponenten mit den Quadratzahlen
fortschreiten, gekannt zu haben scheint«.

Potenzreihen, deren Exponenten eine arithmetische Progression zweiter
Ordnung bilden, treten in \textsc{Jacob Bernoullis} \textit{Ars coniectandi} (Basel, 1713), einem
Werke, das \textsc{Gauss} seit 1793 besessen hat, bei einem Problem der Wahrscheinlichkeitsrechnung
(a. a. O., S. 51—56) auf. Im Jahre 1794 hat \textsc{Gauss} die\par
\clearpage\SourceLeaf{100}\ArithmeticHead{15}{R}{II. DIE ERSTE JUGENDZEIT BIS OKTOBER 1795.}
\noindent \textit{Principia} \textsc{Newtons} erworben\footnote{Es war die Ausgabe Amsterdam 1714; das in der Gaussbibliothek befindliche Buch trägt die Inschrift:
\begin{center}\begin{tabular}{c}C. F. Gauss. 1794\\\hline D. E.\end{tabular}\end{center}
\noindent In einem Briefe an \textsc{Schumacher} vom 12. Mai 1843 (\textit{Briefwechsel Gauss-Schumacher} IV, S. 145) schreibt \textsc{Gauss}: »Mir ist dabei wieder in Erinnerung gekommen, dass ich vor \(\frac12\) Jahrhundert, als ich zuerst \textsc{Newtons} \textit{Principia} las, mehreres unbefriedigend fand .\ .\ .«.}. Wahrscheinlich hat er aus diesem Werke die
Methode der Entwicklung von Funktionen in Potenzreihen kennen gelernt, die
er auch späterhin immer mit grosser Vorliebe und mit besonderem Geschick
gehandhabt und als heuristisches Hilfsmittel bei analytischen Untersuchungen
angewandt hat. — Es ist also sehr wohl möglich, dass er sich damals auch
mit Potenzreihen versucht hat, deren Exponenten eine arithmetische Progression
zweiter Ordnung bilden, und dass ihm ihr Zusammenhang mit dem
agM. aufgefallen war\footnote{Vergl. für diesen Zusammenhang unsern Abriss der Theorie des agM., Nr. 3, Werke X, 1, S. 255.}.

In der historischen Entwicklung sind Reihen, deren Exponenten mit den
Quadratzahlen fortschreiten bezw. arithmetische Progressionen zweiter Ordnung
bilden, schon im XVII. Jahrhundert aufgetreten, nämlich bei der Lösung
des gedachten Problems der Wahrscheinlichkeitsrechnung in Aufsätzen von
\textsc{Jacob Bernoulli} und \textsc{Leibniz}\footnote{Acta eruditorum 1690, S. 219—223 bezw. 358—360.}. \textsc{Euler} hat sich mit der in Rede stehenden
Art von Reihen von 1741 an zu wiederholten Malen beschäftigt\footnote{Siehe \textsc{Euler}, \textit{Observationes analyticae variae de combinationibus}, Comment. Acad. Petrop. 13 (1741—43) 1751, S. 64—93; \textit{Découverte d’une loi tout extraordinaire des nombres par rapport à la somme de leur diviseurs}, Bibliothèque impartiale 3 (1751), S. 10—31, wieder abgedruckt Opera postuma 1, 1862, S. 76—84; \textit{Observatio de summis divisorum}, Novi comment. Acad. Petrop. 5 (1754—55) 1760, S. 59—74. Brieflich hat \textsc{Euler} Entwicklungen dieser Art schon im Januar 1741 an \textsc{Daniel Bernoulli} mitgeteilt, siehe den Brief D. \textsc{Bernoullis} an \textsc{Euler} in P. H. v. \textsc{Fuss}, \textit{Correspondance} II, S. 467, 468, ferner \textsc{Eulers} Briefwechsel mit \textsc{Goldbach} an der Wende der Jahre 1743—44, \textit{Correspondance} I, S. 265, 267, 270, sowie den Brief von \textsc{Nikolaus Bernoulli} an \textsc{Euler} vom 24. Oktober 1742, \textit{Correspondance} II, S. 698. Auch hat \textsc{Daniel Bernoulli} in einem Brief an \textsc{Euler} vom 14. April 1742 die Frage aufgeworfen: »wenn \(a\) ein echter Bruch ist, die Summe der Reihe \(a+a^4+a^9+a^{16}+{}\) etc. zu bestimmen«, \textit{Correspondance} II, S. 493.}. In dem
Briefwechsel, den \textsc{Jacobi} mit P. H. v. \textsc{Fuss} über die Herausgabe der Werke
\textsc{Eulers} geführt hat\footnote{Herausgegeben von \textsc{Stäckel} und \textsc{Ahrens}, Bibliotheca Mathem. (3), 8, 1908, S. 233—306, auch separat erschienen, Leipzig 1908.}, spricht sich \textsc{Jacobi} über die genannten Abhandlungen
\textsc{Eulers} wie folgt aus\footnote{Siehe S. 59, 60 der Separatausgabe.}:\par
\clearpage\SourceLeaf{101}\ArithmeticHead{16}{L}{L. SCHLESINGER, ÜBER GAUSS’ FUNKTIONENTHEORETISCHE ARBEITEN.}
»Sie ist nämlich der erste Fall gewesen, in welchem Reihen aufgetreten
sind, deren Exponenten eine arithmetische Reihe \textls[120]{zweiter} Ordnung bilden,
und auf diese Reihen ist durch mich die Theorie der elliptischen Transzendenten
gegründet worden. Die \textsc{Euler}sche Formel [siehe weiter unten Gl. (4)]
ist ein spezieller Fall einer Formel, welche wohl das wichtigste und fruchtbarste
ist, was ich in reiner Mathematik erfunden habe:
\begingroup\makeatletter\tagsleft@true\let\veqno\leqno\makeatother
\begin{equation*}\tag*{(1)}
\begin{aligned}
&(1-q)(1-q^2)(1-q^3)\ldots(z-z^{-1})\\
&\quad(1-qz^2)(1-q^2z^2)(1-q^3z^2)\ldots\\
&\quad(1-qz^{-2})(1-q^2z^{-2})(1-q^3z^{-2})\ldots\\
&=(z-z^{-1})-q(z^3-z^{-3})+q^3(z^5-z^{-5})\\
&\quad-q^6(z^7-z^{-7})+q^{10}(z^9-z^{-9})\ldots,
\end{aligned}
\end{equation*}\endgroup
wo die Exponenten von \(q\) [d. h.] 1, 3, 6, 10 etc. die \textls[120]{dreieckigen} Zahlen
sind. Setzt man für \(z\) eine imaginäre Kubikwurzel der Einheit, so erhält
man die \textsc{Euler}sche Formel. Hierdurch habe ich sie mit der Trisektion der
elliptischen Integrale in Verbindung gebracht. Differenziert man nach \(z\) und
setzt dann \(z=1\), so erhält man auch für den Kubus des \textsc{Euler}schen Produktes
die schöne Entwicklung
\begingroup\makeatletter\tagsleft@true\let\veqno\leqno\makeatother
\begin{equation*}\tag*{(2)}
\begin{aligned}
&[(1-q)(1-q^2)(1-q^3)\ldots]^3\\
&=1-3q+5q^3-7q^6+9q^{10}-11q^{15}+\text{etc.}\text{«}
\end{aligned}
\end{equation*}\endgroup
\noindent Die Formel (1)\footnote{Die übrigens a. a. O. von \textsc{Jacobi} fehlerhaft wiedergegeben ist, indem auf der linken Seite statt \(z^2\) und \(z^{-2}\) überall \(z\) und \(z^{-1}\) gesetzt ist.} geht, wie \textsc{Jacobi} in den \textit{Fundamenta}\footnote{\textit{Fundamenta nova theoriae functionum ellipticarum} (\textsc{Jacobi} Werke I, S. 234).} zeigt, durch eine
einfache Transformation (man hat \(-qz^2\) für \(z\), dann \(q\) für \(q^2\) zu setzen und
mit \(z\) zu multiplizieren) aus der Formel
\begingroup\makeatletter\tagsleft@true\let\veqno\leqno\makeatother
\begin{equation*}\tag*{(3)}
\begin{aligned}
&(1+qz)(1+q^3z)(1+q^5z)\ldots\left(1+\frac qz\right)\left(1+\frac{q^3}z\right)\left(1+\frac{q^5}z\right)\ldots\\
&=\frac{1+q\left(z+\frac1z\right)+q^4\left(z^2+\frac1{z^2}\right)+q^9\left(z^3+\frac1{z^3}\right)+\ldots}{(1-q^2)(1-q^4)(1-q^6)(1-q^8)\ldots}
\end{aligned}
\end{equation*}\endgroup
hervor, die, wie wir sehen werden, von \textsc{Gauss} 1800 aufgestellt worden ist\footnote{Siehe Werke III, S. 434, letzte Gleichung, Werke X, 1, S. 204, art. [12], ferner aus späterer Zeit Werke III, S. 440, 464.}.
Die von \textsc{Jacobi} erwähnte \textsc{Euler}sche Formel lautet (in den \textsc{Jacobi}schen Zeichen)
\begingroup\makeatletter\tagsleft@true\let\veqno\leqno\makeatother
\begin{equation*}\tag*{(4)}
\begin{aligned}
(1-q)(1-q^2)(1-q^3)\ldots&=1-q-q^2+q^5+q^7-q^{12}-q^{15}+\ldots\\
&=\sum_{\nu=-\infty}^{+\infty}(-1)^\nu q^{\frac{3\nu^2+\nu}{2}},
\end{aligned}
\end{equation*}\endgroup
\clearpage\SourceLeaf{102}\ArithmeticHead{17}{R}{II. DIE ERSTE JUGENDZEIT BIS OKTOBER 1795.}
\noindent sie steht bei \textsc{Euler} ohne Beweis am Schlusse der ersten der oben (Fussnote \textsuperscript{23)})
genannten Abhandlungen\footnote{Vergl. \textit{Briefwechsel Jacobi-Fuss}, Separatausgabe, \textsc{Jacobi} an \textsc{Fuss}, S. 23, \textsc{Fuss} an \textsc{Jacobi}, S. 42; \textit{Fundamenta}, \textsc{Jacobis} Werke I, S. 237 Gl. (6). Bei \textsc{Gauss} findet sich die Formel (4) z. B. Werke III, S. 448, Gl. 23; die letzte von \textsc{Jacobi} angegebene Formel (2) steht bei \textsc{Gauss}, Werke III, S. 440, Zeile 5 v. u. Übrigens hängen diese Identitäten, wie \textsc{Jacobi}, \textit{Fundamenta}, \textsc{Jacobis} Werke I, S. 237 bemerkt, aufs engste mit den von \textsc{Gauss} in der \textit{Summatio quarundam serierum singularium} (1808, Werke II, S. 9ff.) gegebenen Formeln zusammen, vergl. hierzu die Abhandlung \textsc{Jacobis}, \textsc{Crelles} Journal 32, \textsc{Jacobis} Werke VI, S. 302. Weitere historische Notizen über das Auftreten der Formel (4) bei \textsc{Euler} findet man in der letztgenannten \textsc{Jacobi}schen Abhandlung und in dem Aufsatze von P. \textsc{Stäckel}, Bibliotheca mathem. (III) 11, 220ff.}.

Dass \textsc{Gauss} 1794 schon irgend etwas von den \textsc{Euler}schen Entwicklungen
gekannt haben sollte, ist völlig ausgeschlossen; sagt er doch im Vorwort zu
den \textit{Disquisitiones arithm.} (Werke I, S. 6) in Bezug auf seine arithmetischen
Untersuchungen: »als ich zu Anfang des Jahres 1795 mein Interesse dieser
Art von Untersuchungen zuwandte, war ich über alles, was von neueren
Forschern\footnote{Als solche waren vorher genannt \textsc{Fermat}, \textsc{Euler}, \textsc{Lagrange}, \textsc{Legendre}.} auf diesem Gebiete gearbeitet worden war, völlig unwissend und
von allen Hilfsmitteln abgeschnitten, die mir eine Kenntnis hiervon hätten
vermitteln können.« Die Angabe von \textsc{Sartorius}\footnote{A. a. O. \textsuperscript{13)}, S. 15.}, \textsc{Gauss} habe schon auf dem
Carolinum \textsc{Euler} und \textsc{Lagrange} gelesen, steht mit dieser Erklärung von \textsc{Gauss}
in Widerspruch. Die mathematischen Bestände der Bibliotheken von Braunschweig
und Wolfenbüttel beschränkten sich zu jener Zeit auf einige landläufige
Kompendien, und erworben hat \textsc{Gauss} die analytischen Hauptwerke \textsc{Eulers}
(\textit{Introductio, Institutiones calculi differentialis} und \textit{integralis}), nach seiner Mitteilung
an \textsc{Wolfgang Bolyai}\footnote{\textit{Briefwechsel Gauss-Bolyai}, S. 15, Brief vom 30. Dez. 1798.}, erst 1798.

Im übrigen wird man sich vor Augen halten müssen, dass zu jener Zeit
\textsc{Gauss}’ mathematisches Interesse vorwiegend der Zahlentheorie zugewandt war\footnote{Vergl. \textsc{Bachmann}, Abhandlung 1 dieses Bandes, S. 4.};
zu den analytischen Versuchen des jugendlichen \textsc{Gauss} dürfte die Freude am
numerischen Rechnen eine der wesentlichsten Anregungen gegeben haben. —
Diese Freude hat \textsc{Gauss} sich auch späterhin immer bewahrt und sein ganzes
Leben hindurch viel und gern numerisch gerechnet. Er hat nicht nur
bei den geodätischen und astronomischen Arbeiten der späteren Jahre sich
in Bezug auf Genauigkeit der Rechnung nicht genug tun können, sondern\par
\HistoricalFoot{X\textsubscript{2} Abh. 2.}{3}
\clearpage\SourceLeaf{103}\ArithmeticHead{18}{L}{L. SCHLESINGER, ÜBER GAUSS’ FUNKTIONENTHEORETISCHE ARBEITEN.}
\noindent sich, wie wir sehen werden, des numerischen Kalkuls einerseits zur Verifikation
von Resultaten bedient, die er auf analytischem Wege erhalten hatte,
andererseits aber auch dazu, um durch \textls[120]{numerische Induktion} neue Wahrheiten
zu entdecken. Beides hat er mit den grossen Analysten des XVIII.
Jahrhunderts, namentlich mit \textsc{Euler}, gemein. Bei den Nachfolgern von \textsc{Gauss}
verkümmert dieses Verfahren immer mehr und mehr.
